%----------------------------------------------------------------
%
%  File    :  project-guide-libuser.tex
%
%  Author  :  Martin Rabensteiner
% 
%  Created :  31 Feb 2026
% 
%  Changed :  15 May 2026
% 
%----------------------------------------------------------------

\chapter{Library User Guide}
\label{chap:libuser}

This section addresses to those that want to include ListMerger
library in a larger software project. To initialise a page with lists,
the \vname{init.init()} function needs to be called. The parameters
contain an overall HTML id, the ids for the \uiname{Undo} and
\uiname{Redo} buttons, a function that can be called when the
\uiname{Save} button is clicked, an items object, and several strings
for HTML templates. A full example of a configuration can be seen in the
\fname{example/main.js} file.

Items have to be handed over in an object that needs to contain at least
a \vname{merged} key and an \vname{originlists} key. The first can be an
empty array, or contain items. The second one needs to be filled with
the lists as objects with some metadata (\vname{id}, \vname{name}, etc.)
and an \vname{items} key, which also contains an array of items, as seen
in Listing~\ref{list:itemslist}.


\begin{samepage}
\begin{lstlisting}[%
  float=tp,
  aboveskip=\floatsep,
  belowskip=\floatsep,
  xleftmargin=0cm,              % no extra margins for floats
  xrightmargin=0cm,             % no extra margins for floats
  basicstyle=\footnotesize\ttfamily,
  frame=shadowbox,
  numbers=left,
  language=JavaScript,
  label=list:itemslist,
  caption={[JSON Item Lists]%
Example of an initialising object for items. The items in \vname{merged} are rendered to the \uiname{mergelist}, \vname{originlists} are constructed as lists with the base items.
},
]
{
  "merged": [],
  "originlists": [
    {
      "id": "eval1",
      "name": "List 1",
      "items": [
        {
          "id": "li1it1",
          "title": "Item 1",
          "author": "Martin Rabensteiner",
          "category": "Category 1",
          "images": [
            "images/screenshot1.png",
            "images/screenshot2.png"
          ],
          "description": "This is the first item in the list."
        }
      ]
    }
  ]
}

\end{lstlisting}
\end{samepage}

An item needs to have at least an \vname{id}, this id is also used as id
in the HTML DOM. The other keys are optional and need to be included
depending on whether they are needed in the user-defined templates.
Arrays (like images) are always merged together.

Templates define how the items appear in the lists
(\vname{item_template}), their info modals (\vname{dialog_template}),
and the form to merge two items (\vname{merge_template}). In the merge
form, the \vname{name} attribute needs to be ident with the keys used in
an item. A callback function can be handed over to process the data
further, for example to store it in a database. This function is then
called with an array of moved, and merged items. In the example, the
data just is written to the JavaScript console.

\section{Initialising}
To make the library work, the \cname{init()} function needs to be
called. Several parameters can be used to adjust HTML classnames or ids.

\section{Templating}

\section{Inline Editing}
Fields can be configured to be editable inline. For text-only fields,
the \cname{contenteditable} attribute\parencite{mdn-contenteditable}
should be set to \cname{"plaintext-only"}. Setting it to \cname{"true"}
is also sufficient for the functionality, but can lead to unexpected
behaviour when dragging an image over the field.

To use a \elname{select} element for editing, the CSS class
\cname{selectedit} can be put on a \elname{span} element to replace its
content. The \elname{span} needs to have some \cname{data-*} attributes
to build the select box. These attributes are:
\begin{itemize}
  \item \cname{data-name}: Is used as \cname{name} for the
    \elname{select} element and for storing to the item, therefore it
    should be one of the already set children of the item.

  \item \cname{data-options}: Needs to have the possible options of the
    select box. It can be either a one-dimensional array, or an object
    with key-value pairs. If it is an object, the key and values are
    used as they are, othervise the single values are used for both.

  \item \cname{data-values}: Can be a single value, or comma-separted
    the selected values. Valid examples are:
    \begin{itemize}
      \item \lstinline!1,2,3,4!
      \item \lstinline!["John", "Max", "Mario"]!
      \item \lstinline!{"0": "Very Low", "1": "Low", "2": "Medium", "3": "High", "4": "Very High"}!
    \end{itemize}
  
  \item \cname{data-mode}: Per default, the select box is displayed as
    single select. When the mode is \cname{multiple}, the
    \cname{multiple} attribute is added to the \elname{select} element
    as well. Another option for this attribute is \cname{class}, to
    change the class of this element to a specific mode, for example to
    represent colors or icons of the current value.
  
  \item \cname{data-class}: If the mode is set to \cname{class}, this
    value is used as CSS class prefix. When the selection is changed,
    the class consisting of the prefix and the selected value is added
    to this element. Other classes with this prefix are remove on an
    input change.
\end{itemize}

\section{Styling}
Some stylings are highly recommended in order to 
\fname{listmerger.css}
Further recommended stylings can be found in the examples described in
the next chapter.

\section{Examples}
Within the code, a simple and an extended example of implementations are
delivered and can be found in the \fname{/src/example} folder. The
simple example (\fname{simple.html} and \fname{simple.js}) demonstrate a
minimal setup. The extended example is a showcase of nearly all features
implemented in the library. The related stylesheet (\fname{style.css})
can also be very helpful to build an own sophisticated instance.

\section{Customisation}
The selectors of all HTML class names and ids can be set during the
initilisation. The identifiers, default values, and a short description
can be seen in Table~\ref{}

\begin{table}[tp]
\tablerowstretch
\rowcolors{2}{}{tablerowcolour}
\centering
\begin{footnotesize}
\begin{tabularx}{\linewidth}
{>{\leavevmode\kern-\tabcolsep}lm{1cm}lY<{\kern-\tabcolsep}}
\toprule
\textbf{identifiers} & \textbf{Type} & \textbf{Default value} & \textbf{Description} \\
\midrule
% TODO styling
PREFIX\_MOVED & Prefix for ID & "moved-" & Prefix that is added to the HTML IDs of moved items in the mergelist. \\
PREFIX\_MERGED & Prefix for ID & "merged-" & Prefix that is added to the HTML IDs of merged items in the mergelist. \\
\midrule
CONTAINER & ID & "listmerger" & Contains the mergelist and container for origin lists. \\
MERGE\_LIST\_CONTAINER & ID & "mergelist-container" & Is a wrapper for the mergelist and its indicator. \\
MERGE\_LIST & ID & "mergelist" & Is the mergelist itself. \\
LIST & Class & "list" & Every list of items. \\
\midrule
ITEM & Class & "item" & A single item within a list. \\
ITEM\_DRAGHANDLE\_ACTIVE & Class & "draghandle-active" & The item is currently dragged with the draghandle. \\
ITEM\_DRAG\_ACTIVE & Class & "drag-active" & The item is currently dragged itself. \\
\midrule
TAB\_CONTAINER & Class & "tab-container" & The outer container for the tabs and the tab select box. \\
TAB\_BAR & Class & "tabbar" & The innter container for all tabs. \\
TAB\_NAME & Name attribute & "tabs" & Attribute that is added to the \textbackslash{}elname\{details\} tag of an item. \\
TAB\_COMPACT & Class & "compact" & When there is not enough space for a tab titles, this class is added to the tab bar. \\
TAB\_SELECTOR & ID & "detailsselector" & The tab select box shown in the compact state. \\
\midrule
HOVER\_DROP & Class & "drop" & The effect added to a list or item, when an item is dragged over it. \\
ARRANGEBAR & Class & "arrangebar" & The bar shown between two items in the mergelist when rearranging or dragging. \\
\midrule
BUTTON\_MOVE & Class & "move" & The button to move one item to the mergelist or back from there to the origin list. \\
BUTTON\_MOVEALL & Class & "moveall" & The button to move all remaining items of one origin list to the mergelist. \\
BUTTON\_SAVE & Class & "save" & Save button for the edit dialog (currently not in use and replaced with inline editing). \\
BUTTON\_CLOSE & Class & "close" & The Button that cancels the merge dialog. \\
BUTTON\_MERGE & Class & "merge" & The Button that performs the merge after the dialog. \\
BUTTON\_DETACH & Class & "detach" & Added to merged items and their origin items to unpair the merge. \\
BUTTON\_DRAGHANDLE & Class & "draghandle" & Identifies the element of the drag handle template. \\
\midrule
BUTTON\_UNDO & ID & "undo" & The undo button for the history. \\
BUTTON\_REDO & ID & "redo" & The redo button for the history. \\
\midrule
DIALOG\_MERGE\_CONTAINER & Class & "mergecontainer" & Container within the merge dialog. \\
\midrule
TEXT\_BUTTON\_MOVEALL & Text & "Move Remaining" & Shown at every tab. \\
TEXT\_BUTTON\_MERGE & Text & "Merge" & Shown in the merge dialog. \\
TEXT\_BUTTON\_CANCEL & Text & "Cancel" & Shown in the merge dialog. \\
\bottomrule
\end{tabularx}
\end{footnotesize}
\caption[HTML Identifiers]%
{Identifiers used for HTML, CSS, and JavaScript and their description. They can be set while initialising, or otherwise the default values are used.}
\label{tab:identifiers}
\end{table}




\section{Saving}
The current state of the list 

